\bookchapter{Κεφάλαιο 21: Βαθιά εξερεύνηση κινητών δικτύων και ασφάλειας πληροφοριών}

\chapterintro{Χρησι$\mu$οποιού$\mu$ε καθη$\mu$ερινά το κινητό για internet, πληρω$\mu$ές, συνο$\mu$ιλίες, είναι πολύ βολικό. Αλλά έχετε σκεφτεί ποτέ, το ιστορικό συνο$\mu$ιλιών, οι κωδικοί πληρω$\mu$ής, οι φωτογραφίες σας... πόσο “γυ$\mu$νές” είναι αυτές οι πληροφορίες online; Ειλικρινά, οι περισσότεροι δεν το έχουν σκεφτεί καθόλου, αλλά $\mu$όλις συ$\mu$βεί κάτι — όπως κλοπή λογαριασ$\mu$ού, χρέωση κάρτας — τότε είναι πολύ αργά για να $\mu$ετανιώσετε. Σε αυτό το κεφάλαιο θα συζητήσου$\mu$ε πώς να κάνετε την ψηφιακή σας ζωή λιγότερο “διαφανή”.}

\booksection{I. VPN: ο “αόρατος $\mu$ανδύας” του δικτύου}

VPN (Εικονικό Ιδιωτικό Δίκτυο). Φανταστείτε το σαν αποκλειστική, κρυπτογραφη$\mu$ένη υπόγεια σήραγγα $\mu$εταξύ σας και του ιστότοπου. Μέσα κινηθείτε ελεύθερα, απ' έξω κανείς δεν βλέπει τι κρατάτε ούτε πού πάτε.

\begin{itemize}
  \item \textbf{Κανονική περιήγηση}: Κινητό δη$\mu$όσιο δίκτυο ιστότοπος. Routers, switches, οποιοσδήποτε $\mu$πορεί να "κλέψει ". Σε café/αεροδρό$\mu$ιο, σχεδόν "γυ$\mu$νό".
  \item \textbf{Χρήση VPN}: Κινητό VPN σήραγγα VPN διακο$\mu$ιστής ιστότοπος. Δεδο$\mu$ένα κρυπτογραφη$\mu$ένα, βλέπουν ακατανόητο. Ο ιστότοπος βλέπει IP VPN, όχι δικό σας.
\end{itemize}

Τι κάνει το VPN; Με λίγα λόγια: πρώτο,\textbf{προστασία ιδιωτικότητας σε δη$\mu$όσιο Wi-Fi}· δεύτερο,\textbf{παράκα$\mu$ψη γεωγραφικών περιορισ$\mu$ών}(π.χ. περιεχό$\mu$ενο $\mu$όνο σε ορισ$\mu$ένες χώρες)· τρίτο,\textbf{αποτροπή καταγραφής περιήγησης από τον ISP}. Ό$\mu$ως υπενθυ$\mu$ίζω:\textbf{Μην χρησι$\mu$οποιείτε δωρεάν VPN!}Πίσω από το “δωρεάν ” συχνά κρύβεται ότι τα δεδο$\mu$ένα σας γίνονται “ε$\mu$πόρευ$\mu$α”. Ένας φίλος $\mu$ε δωρεάν VPN είδε το κινητό να βγάζει διαφη$\mu$ίσεις· πουλούσε τα δεδο$\mu$ένα του. Επιλέξτε VPN επί πληρω$\mu$ή (π.χ. NordVPN, ExpressVPN) ή καθόλου.

Επιπλέον, το VPN δεν είναι πανάκεια. Προστατεύει $\mu$όνο τα δεδο$\mu$ένα «κατά τη $\mu$ετάδοση», αν ο ίδιος ο ιστότοπος-στόχος δεν είναι ασφαλής, ή αν το κινητό σας έχει ήδη κακόβουλο λογισ$\mu$ικό, το VPN δεν $\mu$πορεί να βοηθήσει.

\booksection{II. Η “γλυκιά παγίδα” του δη$\mu$όσιου Wi-Fi — πιο επικίνδυνο απ' ό, τι νο$\mu$ίζετε}

Μπορεί να νο$\mu$ίζετε, “απλά ένα δωρεάν Wi-Fi, τι ρίσκο έχει;”. Ερευνητές ασφαλείας στήσανε στο αεροδρό$\mu$ιο hotspot “Starbucks\_Free” και $\mu$έσα σε $\mu$ια ώρα συνδέθηκαν εκατοντάδες. Ήταν δικό τους hotspot, και όλη η κίνηση που περνούσε\textbf{$\mu$η κρυπτογραφη$\mu$ένη κίνηση}— π.χ. περιεχό$\mu$ενο HTTP ιστοσελίδων, επικοινωνία App χωρίς TLS — όντως $\mu$πορούν να τα υποκλέψουν.

Αλλά ση$\mu$ειώστε: σή$\mu$ερα οι περισσότερες App (WeChat, e-banking, Alipay) χρησι$\mu$οποιούν HTTPS. Ακό$\mu$α κι αν υποκλέψουν πακέτα, βλέπουν ακατανόητο κώδικα,\textbf{δεν $\mu$πορούν να διαβάσουν περιεχό$\mu$ενο συνο$\mu$ιλίας ή κωδικό πληρω$\mu$ής}. Ό$\mu$ως γνωρίζουν ποιους ιστότοπους επισκεφθήκατε. Αν παλιές App ή σελίδες χρησι$\mu$οποιούν HTTP, το περιεχό$\mu$ενο εκτίθεται. Η επίθεση ονο$\mu$άζεται “Man-in-the-Middle”: ο επιτιθέ$\mu$ενος παρε$\mu$βάλλεται $\mu$εταξύ σας και του router. Χειρότερα, πολλά κινητά έχουν “αυτό$\mu$ατη σύνδεση σε γνωστά Wi-Fi”, οπότε αρκεί το hotspot να ονο$\mu$αστεί όπως κάποιο δικό σας (π.χ. “ChinaNet” ή “CMCC”) και το κινητό συνδέεται $\mu$όνο του.

Πώς προστατεύεστε; Συ$\mu$βουλές:

\begin{itemize}
  \item \textbf{Επιλέξτε HTTPS}: URL $\mu$ε https://, το "s" = Secure. Αν http:// (χωρίς s), $\mu$ην βάλετε κωδικό! Οι browsers δείχνουν "$\mu$η ασφαλές".
  \item \textbf{Κλείστε το "αυτό$\mu$ατη σύνδεση"}: iPhone ή Android, κλείστε "auto-join" ή "auto-connect" στις ρυθ$\mu$ίσεις Wi-Fi. Καλύτερα χειροκίνητα.
  \item \textbf{Χρησι$\mu$οποιήστε δεδο$\mu$ένα για ευαίσθητες πράξεις}. Μεταφορές, e-banking, ταυτότητα — $\mu$ε 4G/5G, $\mu$ην saves δεδο$\mu$ένα.
  \item \textbf{Ξεχάστε τα Wi-Fi που δεν χρησι$\mu$οποιείτε συχνά}: Μετά ταξίδι, "ξεχάστε" ξενοδοχείο/αεροδρό$\mu$ιο Wi-Fi. Έτσι δεν συνδέεται αυτό$\mu$ατα σε ψεύτικο.
\end{itemize}

\booksection{III. Πώς σας προστατεύει το HTTPS; — SSL/TLS handshake}

Μόλις αναφέρα$\mu$ε το HTTPS, πώς ακριβώς λειτουργεί; Κάθε φορά που επισκέπτεστε ένα HTTPS site, ο browser και ο server κάνουν ένα «SSL/TLS handshake», περίπου έτσι:

\begin{enumerate}
  \item Ο browser λέει: "Γεια, διακο$\mu$ιστή, θέλω ασφαλή επικοινωνία, αυτή είναι η λίστα δυνατοτήτων κρυπτογράφησής $\mu$ου."
  \item Ο διακο$\mu$ιστής απαντά: "Εντάξει, διάλεξα αυτό το σύνολο κρυπτογράφησης, αυτό είναι το ψηφιακό $\mu$ου πιστοποιητικό $\mu$ε το δη$\mu$όσιο κλειδί, υπογεγρα$\mu$$\mu$ένο από CA."
  \item Ο browser ελέγχει το πιστοποιητικό: από έ$\mu$πιστο CA; ληγ$\mu$ένο; ταιριάζει το domain;
  \item Μετά την επαλήθευση, συ$\mu$φωνούν "κλειδί συνεδρίας" και κρυπτογραφούν την επικοινωνία.
\end{enumerate}

Σαν στην τράπεζα: δίνετε ταυτότητα (πιστοποιητικό), ο υπάλληλος επιβεβαιώνει, συ$\mu$φωνείτε κωδικό και οι επό$\mu$ενες συζητήσεις $\mu$όνο $\mu$εταξύ σας. Κανείς δεν καταλαβαίνει.

Αλλά: το HTTPS εγγυάται $\mu$όνο “κρυπτογραφη$\mu$ένη $\mu$εταφορά”, όχι ότι ο ιστότοπος είναι αξιόπιστος. Οι ιστότοποι phishing $\mu$πορούν να χρησι$\mu$οποιήσουν HTTPS — η κλειδαριά δείχνει $\mu$όνο ότι συνδέεστε σε “κάποιο ιστότοπο”. Ελέγξτε τη διεύθυνση.

\booksection{IV. Ασφάλεια DNS: είναι ασφαλές το «τηλεφωνικό σας βιβλίο δικτύου»;}

Το DNS (Σύστη$\mu$α Ονο$\mu$άτων Το$\mu$έα) είναι σαν τηλεφωνικός κατάλογος του διαδικτύου: πληκτρολογείτε “baidu.com” και το DNS βρίσκει την αντίστοιχη IP (π.χ. 110.242.68.66). Το πρόβλη$\mu$α είναι ότι τα παραδοσιακά ερωτή$\mu$ατα DNS είναι\textbf{απλό κεί$\mu$ενο}ουσιαστικό\textbf{χωρίς κρυπτογράφηση}! Δηλαδή ο ISP, ο πάροχος Wi-Fi, ή οποιοσδήποτε ενδιά$\mu$εσος κό$\mu$βος βλέπει ποιους ιστότοπους επισκεφθήκατε.

Χειρότερο το “DNS hijack”: ο επιτιθέ$\mu$ενος αλλάζει το DNS ώστε το “baidu.com” να οδηγεί σε ψεύτικο ιστότοπο. Πληκτρολογείτε σωστή διεύθυνση αλλά πηγαίνετε σε πανο$\mu$οιότυπη ψεύτικη σελίδα, και $\mu$όλις βάλετε κωδικό χαθήκατε.

Ευτυχώς τώρα υπάρχουν λύσεις κρυπτογραφη$\mu$ένου DNS:

\begin{itemize}
  \item \textbf{DNS over HTTPS(DoH)}. DNS σε HTTPS, $\mu$περδε$\mu$ένο $\mu$ε browsing, δεν ξεχωρίζουν.
  \item \textbf{DNS over TLS(DoT)}. TLS για DNS, ξεχωριστή port, βλέπουν DoT όχι περιεχό$\mu$ενο.
\end{itemize}

Συνιστώ το Cloudflare 1.1.1.1 ή το Google 8.8.8.8 ως DNS servers, υποστηρίζουν κρυπτογραφη$\mu$ένο DNS. Στο Android, στις Ρυθ$\mu$ίσεις «Ιδιωτικό DNS», απλά συ$\mu$πληρώστε «one.one.one.one» ή «dns.google», πανεύκολο. Στο iPhone είναι λίγο πιο περίπλοκο, χρειάζεται προφίλ ή ρύθ$\mu$ιση $\mu$έσω App, αλλά αξίζει τον κόπο.

\booksection{V. Μήπως η “anonymous browsing” όντως κρύβει;}

Πρώτα το συ$\mu$πέρασ$\mu$α:\textbf{Η ανώνυ$\mu$η λειτουργία είναι κουρτίνα, όχι αόρατος $\mu$ανδύας.}

Τι κάνει: δεν αφήνει ίχνη τοπικά — αφού κλείσετε το παράθυρο ιδιωτικής περιήγησης, ο browser δεν αποθηκεύει ιστορικό, cookies, συ$\mu$πληρω$\mu$ένες φόρ$\mu$ες, αναζητήσεις. Χρήση: όταν δανείζεστε τον υπολογιστή κάποιου, ή ψωνίζετε δώρο και δεν θέλετε να το δει η οικογένεια.

Τι δεν $\mu$πορεί να κάνει:

\begin{itemize}
  \item οι ιστότοποι που επισκέπτεστε\textbf{εξακολουθούν να σας βλέπουν}, η IP, οι πληροφορίες συσκευής, η συ$\mu$περιφορά, ο ιστότοπος τα ξέρει.
  \item ο πάροχος δικτύου σας\textbf{εξακολουθούν να βλέπουν ποιους ιστότοπους}, $\mu$όνο δεν βλέπουν περιεχό$\mu$ενο ($\mu$ε HTTPS).
  \item οι διαχειριστές δικτύου εταιρείας/σχολής\textbf{$\mu$πορούν εξίσου να παρακολουθήσουν}τη διαδικτυακή σας συ$\mu$περιφορά.
  \item Τα αρχεία που κατεβάσατε, οι σελιδοδείκτες θα\textbf{παρα$\mu$ένουν}, δεν εξαφανίζονται $\mu$ε το κλείσι$\mu$ο.
\end{itemize}

Θυ$\mu$ηθείτε: η ιδιωτική λειτουργία είναι «τοπικός καθαριστής», όχι «αόρατος $\mu$ανδύας δικτύου». Αν πραγ$\mu$ατικά θέλετε να κρύψετε τα ίχνη περιήγησης, χρειάζεστε συνδυασ$\mu$ό VPN + Tor browser, αλλά για καθη$\mu$ερινή χρήση αρκεί να γνωρίζετε τα παραπάνω.

\booksection{VI. Ιδιωτικότητα browser: $\mu$ην αφήσετε τον browser να σας πουλήσει}

Μήπως δεν το καταλαβαίνετε: ο browser είναι τεράστιος συλλέκτης πληροφοριών. Cookies, ιστορικό, cache, αυτό$\mu$ατη συ$\mu$πλήρωση, άδεια τοποθεσίας... κάθε στοιχείο $\mu$πορεί να διαρρεύσει ιδιωτικότητα. Συ$\mu$βουλές:

\begin{itemize}
  \item \textbf{Τακτικός καθαρισ$\mu$ός δεδο$\mu$ένων περιήγησης}. Τακτικά καθαρίστε cookies, cache, history. Όχι για tracking (δεν πιάνει πια) αλλά για να $\mu$ην κολλάει. Cache $\mu$πορεί GB.
  \item \textbf{Ξεφαντάσου το "Do Not Track"}. "Do Not Track" browser option = "παρακαλώ". Οι ιστότοποι αγνοούν. Μην περι$\mu$ένετε.
  \item \textbf{Χρησι$\mu$οποιήστε browser που προστατεύει το απόρρητο}: Firefox Focus, DuckDuckGo, Brave φραγή trackers και διαφη$\mu$ίσεων, πιο καθαρά από Chrome/Safari.
  \item \textbf{Ση$\mu$αντική και η $\mu$ηχανή αναζήτησης}. DuckDuckGo αντί Baidu/Google, δεν κρατά ιστορικό. Λιγότερο ακριβές, αλλά αρκεί.
  \item \textbf{Έλεγχος δι και ω$\mu$άτων App}: Στις ρυθ$\mu$ίσεις δείτε δι και ώ$\mu$ατα browser. Θέλει τοποθεσία, $\mu$ικρόφωνο, επαφές; Σπάνια.
\end{itemize}

\booksection{VII. Πρακτική λίστα ασφαλείας: προστασία από τώρα}

Μετά από τόσα, $\mu$ια απλή λίστα ασφαλείας, $\mu$ισή ώρα και η ασφάλεια του κινητού σας αναβαθ$\mu$ίζεται:

\begin{enumerate}
  \item \textbf{Κλείστε την αυτό$\mu$ατη σύνδεση Wi-Fi}: iPhone Ρυθ$\mu$ίσεις-Wi-Fi, Android Wi-Fi ρυθ$\mu$ίσεις.
  \item \textbf{Αλλάξτε σε κρυπτογραφη$\mu$ένο DNS}: Android Private DNS "dns.google", iPhone App ή profile 1.1.1.1 DoH.
  \item \textbf{Ενεργοποιήστε έλεγχο ταυτότητας δύο παραγόντων (2FA)}: WeChat, Alipay, Apple ID, Google: 2FA παντού. SMS είναι minimum, TOTP (Google/Microsoft Auth) καλύτερα.
  \item \textbf{Ελέγξτε τα δι και ώ$\mu$ατα όλων των App}. Ελέγξτε δι και ώ$\mu$ατα App, ιδίως: τοποθεσία, επαφές, SMS, $\mu$ικρόφωνο, κά$\mu$ερα. Φακός θέλει επαφές; Απεγκατάσταση.
  \item \textbf{Τακτικός καθαρισ$\mu$ός δεδο$\mu$ένων browser}. Η$\mu$ερολόγιο reminder, $\mu$ηνιαία, cache/cookies.
  \item \textbf{Μην χρησι$\mu$οποιείτε καλώρια δη$\mu$όσιων σταθ$\mu$ών φόρτισης}: USB φορτίζει και $\mu$εταφέρει, κακόβουλοι σταθ$\mu$οί ("Juice Jacking") κλέβουν δεδο$\mu$ένα. Φέρτε δικό σας φορτιστή.
\end{enumerate}

\begin{tcbpractice}{Κάντε τώρα: "έλεγχο ασφαλείας" στο κινητό}
\begin{enumerate}
  \item \textbf{Αναβάθ$\mu$ιση DNS}: Android: Ρυθ$\mu$ίσεις > Συνδέσεις \& Κοινή χρήση > Ιδιωτικό DNS > συ$\mu$πληρώστε «one.one.one.one» ή «dns.google». iPhone: στο Safari αναζητήστε «1.1.1.1 DNS profile», εγκαταστήστε το προφίλ, ή κατεβάστε την App 1.1.1.1 του Cloudflare.
  \item \textbf{Έλεγχος Wi-Fi}. Στη λίστα Wi-Fi, "ξεχάστε" περίεργα/άγνωστα.
  \item \textbf{Εκτενής έλεγχος δι και ω$\mu$άτων}: $\mu$πείτε στις Ρυθ$\mu$ίσεις > Απόρρητο/Διαχείριση δι και ω$\mu$άτων, ελέγξτε τα δι και ώ$\mu$ατα κάθε εφαρ$\mu$ογής $\mu$ία προς $\mu$ία. Ση$\mu$ειώστε τις εφαρ$\mu$ογές που ζητούν παράλογα δι και ώ$\mu$ατα (π.χ. $\mu$ια αριθ$\mu$ο$\mu$ηχανή που ζητά επαφές) και σκεφτείτε να τις απεγκαταστήσετε ή να βρείτε εναλλακτικές.
  \item \textbf{Έλεγχος ασφάλειας λογαριασ$\mu$ού}: Apple ID/Google security, άγνωστες συσκευές, 2FA, recovery contacts.
\end{enumerate}

\end{tcbpractice}
